%!TEX root = ../../../standalone/applications/app-eki.tex
\ApplicationFicheHeading{app-eki}
\textbf{app-eki: Ensemble Kalman Inversion}
\label{sec:app:specs:app-eki}

% Generated by scripts/sync_applications.py; edit the CSV snapshot.
\begin{FicheMetadataTable}
id & app-eki \\
name & EKI \\
Partners & Inria BXSO, U Grenoble Alpes \\
PC & PC1 - ExaMA \\
Responsible (Permanent) & M. Asch \\
WP7 Engineer & \todo[app-eki / WP7 Engineer]{Sheet field missing.} \\
work\_package & WP4 \\
application\_type & mini-app \\
purpose & Demonstrate feasibility of ensemble Kalman filter based inversion of a simplified wave equation \\
Method-Algorithm WP1 & spectral element, dG/hdG \\
Method-Algorithm WP4 & stochastic inverse problem, stochastic data assimilation: ensemble \\
Method-Algorithm WP6 & uncertainty propagation \\
WP7 & Geophysics \\
inputs & Gmsh, in-house \\
outputs & in-house \\
metrics & manufactured-solution, weak-scalability \\
status & spec-only \\
Benchmark scope & ensemble-mode \\
Framework & PETSc, GeoS \\
parallel\_framework & MPI \\
repo\_url & \todo[app-eki / repo\_url]{Sheet field missing.} \\
tex\_url & \todo[app-eki / tex\_url]{Sheet field missing.} \\
\end{FicheMetadataTable}


\subsection{Description of the benchmark}
\textbf{Status:} Proposed (spec-only) | \textbf{Partners:} Inria BXSO, U Grenoble Alpes | \textbf{Target:} 1/15/2027

\textbf{Purpose:} Demonstrate feasibility of ensemble Kalman filter based inversion of a simplified wave equation for geophysics applications.

\textbf{Key Methods:}
\begin{itemize}
\item WP1: Spectral element, dG/hdG discretization
\item WP4: Stochastic inverse problems, ensemble-based data assimilation
\item WP6: Uncertainty propagation
\end{itemize}

\textbf{Framework:} PETSc, GeoS | \textbf{Parallel:} MPI (ensemble mode)

\textbf{Metrics:} manufactured-solution validation, weak-scalability

\textbf{Scientific Context:} Ensemble Kalman inversion provides a derivative-free approach to inverse problems by treating inversion as a filtering problem. The method is particularly attractive for high-dimensional parameter spaces where gradient computation is expensive or unavailable.

\subsection{Benchmarking methodology}
\begin{description}
  \item[Method-Algorithm WP1:] spectral element, dG/hdG
  \item[Method-Algorithm WP4:] stochastic inverse problem, stochastic data assimilation: ensemble
  \item[Method-Algorithm WP6:] uncertainty propagation
  \item[Framework:] PETSc, GeoS
  \item[parallel\_framework:] MPI
  \item[Benchmark scope:] ensemble-mode
  \item[metrics:] manufactured-solution, weak-scalability
\end{description}
\todo[app-eki / benchmarking methodology]{Benchmarking methodology to be supplied by contributors.}

\subsection{Input/Output Dataset Description}

\subsubsection{Input Data}
Gmsh, in-house

\subsubsection{Output Data}
in-house

\subsection{Results summary}
\todo[app-eki / results summary]{Results summary to be supplied by contributors.}
